\documentclass[11pt]{article}
\usepackage[T1]{fontenc}
\usepackage[utf8]{inputenc}
\usepackage{lmodern,amsmath,amssymb,amsthm,mathtools,microtype,booktabs,longtable}
\usepackage[margin=27mm]{geometry}
\usepackage[round,authoryear]{natbib}
\usepackage{xurl}
\usepackage[hidelinks]{hyperref}
\newcommand{\workbenchurl}{https://github.com/KokunoYumeto/arithmetic-geometry-complexity-workbench}
\newcommand{\recordurl}{https://github.com/KokunoYumeto/mathematics-workbenches}
\newcommand{\C}{\mathbb C}
\newcommand{\R}{\mathbb R}
\newcommand{\Q}{\mathbb Q}
\newcommand{\Z}{\mathbb Z}
\newcommand{\LS}{\operatorname{LS}}
\DeclareMathOperator{\per}{per}
\DeclareMathOperator{\tr}{tr}
\DeclareMathOperator{\St}{St}
\newtheorem{theorem}{Theorem}[section]
\newtheorem{proposition}[theorem]{Proposition}
\newtheorem{lemma}[theorem]{Lemma}
\newtheorem{corollary}[theorem]{Corollary}
\theoremstyle{definition}
\newtheorem{definition}[theorem]{Definition}
\theoremstyle{remark}
\newtheorem{remark}[theorem]{Remark}
\title{Level spaces of the permanent at the all-ones matrix}
\author{}
\date{25 November 2026}
\begin{document}
\maketitle
\begin{abstract}
Minc's expansion expresses the permanent of the all-ones matrix plus a perturbation through its weighted rook numbers. We study linear spaces on which every nonconstant term vanishes, a question arising in the public P versus NP programme and its lower bounds for border determinantal complexity. Using Mirandola and Z\'emor's coordinatewise-product theorem, we prove the unconditional bound $(n^2+2n-1)/3$ in every order at least three; in particular the order-six upper bound becomes fifteen. For order six we also prove the bound ten for every level space invariant under a full row or column cycle, and for every space on which a nontrivial row/column symmetry acts as a scalar. We also identify an eight-by-eight Klein slice exactly. We give the connected-partition expansion of the logarithm of the rook polynomial, including every term in degrees five and six, and prove dimension bounds, a Taylor-order obstruction, amplification, and a classification of rank-one level lines. These results delimit the problem; they do not establish a uniform linear bound or the proposed bound ten for matrices of order six. The mathematics and text were produced by GPT-6 Astra Pro, with exact rational checks; separate re-proofs and their scope are recorded with the paper, and no human has checked it.
\end{abstract}

\section{The question and the results}
The question is how large a linear family of perturbations can leave a permanent exactly constant.
For an $n\times n$ matrix $A$, write
\[
 \per(A)=\sum_{\sigma\in S_n}\prod_{i=1}^n A_{i,\sigma(i)},
 \qquad J_n=(1)_{1\leq i,j\leq n}.
\]
A complex linear subspace $L\subseteq\C^{n\times n}$ is a \emph{level space} if
$\per(J_n+t)=n!$ for every $t\in L$. Let $\LS(n)$ be the largest possible dimension.
The public P versus NP programme asks whether $\LS(n)\leq Cn$ for a constant independent of $n$ \citep{Fleet}.
The first numerical target considered here is $\LS(6)\leq10$.
Theorem~\ref{thm:weighted-bound} proves the unconditional bound $\LS(6)\le15$. Theorem~\ref{thm:rowcycle} proves the target under full row- or column-cycle invariance; Theorem~\ref{thm:scalar-six} proves it under a nontrivial scalar permutation symmetry. Theorem~\ref{thm:hook} shows that no larger level space contains the standard hook. These are restrictions on possible counterexamples, not reductions of arbitrary spaces to symmetric ones. The elementary construction below gives $\LS(6)\geq9$.
Neither the target nor the uniform linear bound is proved in this paper.

The geometric reason for studying the problem is a transfer theorem in the public proof repository: a bound $\LS(n)\leq G$ implies that every border determinantal representation of the permanent of order $n+1$, of size $N$, satisfies
\begin{equation}\label{eq:transfer}
 (n+1)^2+1\leq G+2(n+1)+N.
\end{equation}
The exact Lean statement is the theorem on level-space transfer in the module cited in \citet{FleetTransfer}; it specifies the border representation hypothesis.
For $n=6$ and $G=10$, the inequality is $50\leq24+N$, hence $N\geq26$.
This is a conditional consequence of the target ten, which is not established here. The unconditional cubic theorem does give the following bound.

\begin{corollary}[Consequence for border determinantal size]\label{cor:border-cubic-bound}
For $m\ge4$, every coefficientwise limit of affine determinant expressions of size $N$ for the permanent of order $m$ satisfies
\[
 N\ge m^2-2m+1-\left\lfloor\frac{m^2-2}{3}\right\rfloor.
\]
\end{corollary}
\begin{proof}
Apply the cited level-space transfer with $n=m-1$ and
$G=\lfloor((m-1)^2+2(m-1)-1)/3\rfloor=\lfloor(m^2-2)/3\rfloor$, using Theorem~\ref{thm:weighted-bound}. Rearranging \eqref{eq:transfer} gives the assertion.
\end{proof}
For example the formula gives $N\ge74$ at $m=12$, compared with the classical border lower bound $m^2/2$ of \citet[Theorem~1.0.1, p.~1 of the April~2010 author manuscript]{LMR}. This is a finite numerical consequence of the level-space theorem; it is not an assertion that this asymptotic bound is the strongest available.
Mignon and Ressayre's quadratic lower bound concerns determinantal complexity \citep[Theorem~1]{MR}; the distinction between an exact representation and its closure must be kept when comparing results.

For a rectangular matrix $t$, let $S_k(t)$ be the sum of its $k\times k$ subpermanents, with $S_0=1$ and $S_k=0$ above the smaller side length. Define
\[
 R_t(z)=\sum_{k\geq0}S_k(t)z^k,
 \qquad c_k(t)=[z^k]\log R_t(z).
\]
The logarithm is a formal power series over $\Q$; no analytic convergence is involved. A matrix is \emph{doubly centred} if every row sum and every column sum is zero.

\begin{theorem}\label{thm:summary}
For complex square matrices the following statements hold.
\begin{enumerate}
\item A linear space is level precisely when all its positive-degree rook numbers vanish. Equivalently, its first $n$ rook cumulants vanish.
\item The cumulants have the connected-partition formula of Theorem~\ref{thm:connected}. On doubly centred matrices the degree-five and degree-six expansions are exactly those in Table~\ref{tab:cumulants}.
\item For $n\geq2$,
\[
 2n-3\leq\LS(n)\leq
 \min\left\{\left\lfloor\frac{n^2-1}{2}\right\rfloor,
 \left\lfloor\frac{n^2+2n-1}{3}\right\rfloor\right\}.
\]
A real level space has dimension at most $2n-2$.
\item The function $\LS$ is superadditive, and
\[
 \lim_{n\to\infty}\frac{\LS(n)}n
 =\sup_{n\geq1}\frac{\LS(n)}n,
\]
where infinity is permitted.
\end{enumerate}
\end{theorem}

In words, the first three homogeneous equations give an unconditional quadratic bound. The improvement uses the marginal sums directly, instead of first paying their full dimension and bounding the doubly centred kernel. The real quadratic form also forces a real level space to be much smaller than a general complex isotropic space. They do not supply the higher-degree dimension inequality that a linear bound requires. The Taylor-order theorem below makes that limitation quantitative.


\begin{theorem}[Symmetric spaces and the Klein slice]\label{thm:symmetric-summary}
Every complex level space of six-by-six matrices invariant under a full row or column six-cycle has dimension at most ten. The same bound holds if a nonidentity scalar $\omega\ne1$ is the action on the whole space of a row permutation, column permutation, or a composition of these with transposition. For every $p,q\ge2$, the standard level space supported in the first row and column, of dimension $p+q-3$, admits no proper linear level-space extension.

Let $\operatorname{PSL}_2(\mathbb F_7)$ act simultaneously on the rows and columns indexed by $\mathbb P^1(\mathbb F_7)$. Each of its two conjugate non-real three-dimensional matrix summands has exactly twenty-four level lines, the cone over the ordinary flexes of Klein's quartic. Their direct sum contains no level space of dimension three, and a level space containing one of these pure lines has dimension one. The existence of a level plane transverse to both summands is not decided.
\end{theorem}
The symmetry conclusions are proved in Theorems~\ref{thm:rowcycle}, \ref{thm:scalar-six} and \ref{thm:hook}. The exact Klein equations, including their normalization over $\mathbb Q(\sqrt{-7})$, are Theorem~\ref{thm:klein-pure} and Propositions~\ref{prop:klein-mixed}--\ref{prop:klein-noextension}. These conclusions exclude specified classes of larger spaces and make no claim that every extremal space belongs to one of those classes.

\subsection{Sources and scope}
The permanent expansion is classical; the public proof repository cites \citet{Minc} and gives its exact formal statement \citep{FleetExpansion}. Partition-lattice inversion for rook polynomials goes back to \citet{GJW}; the connected formula used here is its labelled exponential-formula refinement. We provide the refinement and the small-degree calculations in full. Our quadratic and cubic dimension argument uses the coordinatewise-product theorem of \citet[Theorem~3.3, preprint p.~7]{MZ}.

The Fano schemes studied by \citet{CI} parametrize linear spaces on homogeneous determinantal and permanental varieties. Their diagonal-torus argument, in Section~2.2 and Proposition~2.3, does not preserve the translated level condition at $J_n$. Section~\ref{sec:symmetry} explains this obstruction explicitly.
OpenAI's September preprint states a cubic border determinantal lower bound for arbitrary affine expressions and coefficientwise limits, with explicit constants \citep[Theorem~1.1, p.~2]{OpenAICubic}. More recent preprints of Anthropic and OpenAI claim superpolynomial arithmetic and border lower bounds \citep{Anthropic,OpenAI}. If valid, these results are asymptotically stronger than the numerical transfer proved here. They do not classify level spaces, and none of the proofs below assumes their claims. Their theorem statements were consulted; their proofs have not been independently verified here.

\section{Proof overview}
The permanent expansion first replaces one polynomial identity on a subspace by homogeneous identities. We then enforce distinct row and column labels by two applications of partition-lattice inversion. A pair of partitions is a bipartite multigraph; the logarithm keeps its connected components. Double centring kills every graph with a vertex of degree one. This gives finite, checkable expressions for the higher conditions.

The dimension arguments use different parts of that description. The quadratic condition is a nondegenerate form on the total-sum-zero hyperplane. For the unconditional cubic bound, we replace each entry by that entry minus half its row and column sums, then append scaled marginal sums. The quadratic and cubic equations become two moments for a common weighted coordinate form, and the first moment vanishes too. Mirandola and Z\'emor's product-dimension theorem applies after adjoining the constant vector and saturating by the product stabilizer. The sole possible zero-weight block at orders four and six is excluded by a marginal-sum constraint. Root-of-unity row blocks show why a fixed list of such identities cannot yield a uniform linear bound. Block sums amplify every example, while rank-one matrices reduce all the conditions to disjoint supports of two elementary-symmetric generating polynomials.

\section{Rook numbers and connected partitions}
This section obtains formulas valid in every matrix size, rather than formulas inferred by interpolation.

\begin{proposition}[Minc's expansion]\label{prop:minc}
For $t\in\C^{n\times n}$,
\[
 \per(J_n+t)=\sum_{k=0}^n(n-k)!S_k(t).
\]
Consequently a complex linear space $L$ is level if and only if $S_k|_L=0$ for $1\leq k\leq n$.
\end{proposition}
\begin{proof}
In the expansion of $\prod_i(1+t_{i,\sigma(i)})$, choose the $k$ positions at which the perturbation is used. A fixed matching of $k$ rows to $k$ columns extends to exactly $(n-k)!$ permutations. Summing over these matchings gives the formula.
If $L$ is level, then for every $t\in L$ the polynomial $\per(J_n+zt)$ is constant for all complex $z$, since $zt\in L$. Its degree-$k$ coefficient is $(n-k)!S_k(t)$, which must vanish. The converse follows from the expansion.
\end{proof}

Let $\Pi_k$ be the set of partitions of $\{1,\ldots,k\}$. For $\pi\in\Pi_k$, put
\[
 \mu(\pi)=\prod_{B\in\pi}(-1)^{|B|-1}(|B|-1)!.
\]
For $\pi,\sigma\in\Pi_k$, define
\[
 T_{\pi,\sigma}(t)=
 \sum_{a:\pi\to[n]}\sum_{b:\sigma\to[n]}
 \prod_{\ell=1}^k t_{a(B_\pi(\ell)),b(B_\sigma(\ell))}.
\]
The labels in these sums need not be distinct. The two partitions are \emph{connected} when their join is the one-block partition. Equivalently, the bipartite multigraph whose vertices are the blocks, and whose edges are the labels $1,\ldots,k$, is connected.

\begin{theorem}\label{thm:connected}
For every matrix $t$ and every $k\geq1$,
\begin{align}
 S_k(t)&=\frac1{k!}\sum_{\pi,\sigma\in\Pi_k}
 \mu(\pi)\mu(\sigma)T_{\pi,\sigma}(t),\label{eq:mobius}\\
 c_k(t)&=\frac1{k!}\sum_{\substack{\pi,\sigma\in\Pi_k\\\pi\vee\sigma=\hat1}}
 \mu(\pi)\mu(\sigma)T_{\pi,\sigma}(t).\label{eq:connected}
\end{align}
If all row and column sums of $t$ vanish, terms with any singleton block are zero and may be omitted.
\end{theorem}
\begin{proof}
For a tuple of row indices, partition the positions having equal indices. The indicator that all indices are distinct is the sum of $\mu(\pi)$ over partitions whose blocks impose equalities satisfied by that tuple. This is the defining inversion identity for the partition lattice. Apply the same identity independently to column indices. The resulting sum counts ordered pairs of injective tuples, with their matrix-entry products. Each matching is counted $k!$ times, proving \eqref{eq:mobius}.

Every pair of partitions decomposes uniquely into its connected components on disjoint sets of labels. Both the M\"obius factors and the unrestricted-label sum factor over these components: equal numerical row labels in different components impose no restriction. If the sum of weights of connected pairs on $j$ labels is $w_j$, this unique decomposition gives
\[
 \sum_{k\geq0}S_kz^k
 =\sum_{r\geq0}\frac1{r!}\left(\sum_{j\geq1}\frac{w_j}{j!}z^j\right)^r.
\]
For completeness, the coefficient counts a partition of the $k$ labels into $r$ unordered nonempty components, then one connected pair on each component. The factors $r!$ and $j!$ account exactly for their orderings. Taking the formal logarithm proves \eqref{eq:connected}.

If a row-partition block is a singleton, its numerical label occurs in one factor only. Summing over that label gives a column sum of $t$, hence zero. A column singleton similarly gives a row sum.
\end{proof}

For a nonnegative integer matrix $A=(a_{ij})$ of size $r\times s$, write
\[
 T_A(t)=\sum_{u_1,\ldots,u_r=1}^n\sum_{v_1,\ldots,v_s=1}^n
           \prod_{i=1}^r\prod_{j=1}^s t_{u_i v_j}^{a_{ij}},
 \qquad T_A^+=T_A+T_{A^{\mathsf T}}.
\]
Zero exponents contribute one. Rows and columns of $A$ are unlabeled for this notation, since permuting them does not change $T_A$.
The small-degree formulas can now be written compactly.
Let $p_j=\sum_{i,l}t_{il}^j$, $R_4=\sum_i(\sum_l t_{il}^2)^2$, $C_4=\sum_l(\sum_i t_{il}^2)^2$, and $Q_4=\tr((tt^{\mathsf T})^2)$.
On doubly centred matrices,
\begin{align*}
 c_2&=p_2/2,&c_3&=2p_3/3,\\
 c_4&=\tfrac32p_4-\tfrac34(R_4+C_4)+\tfrac14Q_4.
\end{align*}
Since $S_1=0$, this also gives
$S_4=c_4+S_2^2/2$.

\begin{table}[!htbp]
\centering\small
\begin{tabular}{@{}cl@{}}
\toprule
Degree&Connected contraction and coefficient\\\midrule
5&$\displaystyle \frac{24}{5}T_{(5)}-4T_{(2\ 3)}^+
 +T_{\left(\begin{smallmatrix}0&2\\2&1\end{smallmatrix}\right)}
 +2T_{\left(\begin{smallmatrix}1&1\\1&2\end{smallmatrix}\right)}$\\[10pt]
6&$\displaystyle 20T_{(6)}-15T_{(2\ 4)}^+-\frac{20}{3}T_{(3\ 3)}^+
 +\frac52T_{(2\ 2\ 2)}^+$\\[8pt]
&$\displaystyle {}+\frac92T_{\left(\begin{smallmatrix}0&2\\2&2\end{smallmatrix}\right)}
 +2T_{\left(\begin{smallmatrix}0&2\\3&1\end{smallmatrix}\right)}^+
 +6T_{\left(\begin{smallmatrix}1&1\\1&3\end{smallmatrix}\right)}$\\[8pt]
&$\displaystyle {}+3T_{\left(\begin{smallmatrix}1&1\\2&2\end{smallmatrix}\right)}^+
 +2T_{\left(\begin{smallmatrix}1&2\\2&1\end{smallmatrix}\right)}$\\[8pt]
&$\displaystyle {}-\frac32T_{\left(\begin{smallmatrix}0&1&1\\2&1&1\end{smallmatrix}\right)}^+
 -\frac12T_{\left(\begin{smallmatrix}0&1&2\\2&1&0\end{smallmatrix}\right)}^+
 -\frac13T_{\left(\begin{smallmatrix}1&1&1\\1&1&1\end{smallmatrix}\right)}^+$\\[8pt]
&$\displaystyle {}+\frac16T_{\left(\begin{smallmatrix}0&1&1\\1&0&1\\1&1&0\end{smallmatrix}\right)}$\\
\bottomrule
\end{tabular}
\caption{The complete formulas for $c_5$ and $c_6$ on doubly centred matrices. The notation $T_A^+$ includes the transpose even when the numbers of its row and column vertices differ.}\label{tab:cumulants}
\end{table}

\begin{proof}[Verification of the table]
For each bipartite multigraph $A$, let $\operatorname{Aut}_{\mathrm{bip}}(A)$ be the permutations of its row and column vertices preserving $A$. The coefficient in \eqref{eq:connected} is
\begin{equation}\label{eq:graphweight}
 \frac{\prod_i(-1)^{r_i-1}(r_i-1)!
       \prod_j(-1)^{s_j-1}(s_j-1)!}
 {|\operatorname{Aut}_{\mathrm{bip}}(A)|\prod_{i,j}a_{ij}!},
\end{equation}
where $r_i$ and $s_j$ are its degrees. Indeed, the number of labellings of its $k$ edges is $k!$ divided by its parallel-edge factorials and its vertex automorphism group; the $k!$ cancels \eqref{eq:connected}.
All degrees are at least two. For total degree five, the two vertex counts are at most two; listing their possible degree sums yields the five types displayed, with transposes counted separately. For total degree six they are at most three. The possible degree partitions on either side are $(6)$, $(4,2)$, $(3,3)$ and $(2,2,2)$. Listing the connected contingency matrices with each pair of these margins gives the twenty-one types displayed, again including transposes. Formula \eqref{eq:graphweight} gives every coefficient. For example, the doubled-single-edge graph has coefficient $1/2$; the four-cycle has coefficient $1/4$; and the six-cycle has coefficient $1/6$. These checks also fix the normalization.
\end{proof}

\section{Elementary dimension bounds}
This section separates consequences of the first two nontrivial conditions from claims requiring all orders.
On the hyperplane $S_1=0$, an elementary two-matching count gives
\begin{equation}\label{eq:quad}
 2S_2(t)=\sum_{ij}t_{ij}^2-\sum_i\left(\sum_jt_{ij}\right)^2
                         -\sum_j\left(\sum_it_{ij}\right)^2.
\end{equation}
Every matrix there decomposes uniquely as
$t=D+a\mathbf1^{\mathsf T}+\mathbf1b^{\mathsf T}$,
where $D$ is doubly centred and $\sum a_i=\sum b_j=0$.
The three summands are orthogonal in \eqref{eq:quad}; its coefficients relative to the entrywise form are $1,1-n,1-n$. Thus it is nondegenerate for $n\geq2$, with real signature $((n-1)^2,2n-2)$.
Polarization makes a level space totally isotropic. Over $\C$ its dimension is therefore at most $\lfloor(n^2-1)/2\rfloor$; over $\R$ it is at most $\min\{(n-1)^2,2n-2\}$.

For subspaces $U,V\subseteq\C^M$, let $UV$ be the linear span of their coordinatewise products, and write $V^2=VV$. The coordinatewise stabilizer is
\[
 \St(V)=\{a\in\C^M:aV\subseteq V\}.
\]
It is a unital subalgebra of $\C^M$. Orthogonal complements in the next lemma use the bilinear form $\langle x,y\rangle=\sum_i x_i y_i$, without complex conjugation.

\begin{lemma}\label{lem:fermat}
If a subspace $W\subseteq\C^M$ satisfies $\sum_iw_i^2=\sum_iw_i^3=0$ for every $w\in W$, then $3\dim W\leq M$.
\end{lemma}
\begin{proof}
Polarization gives $W+W^2\subseteq W^\perp$, with coordinatewise products. Put $V=W+W^2$ and $H=\St(V)$. Every unital subalgebra of $\C^M$ consists of vectors constant on the parts of a coordinate partition: this follows by separating any two coordinates not identified by all its elements and using polynomial interpolation to construct the part indicators. The indicators of $H$ decompose $V$ into coordinate blocks $V_B$, each with scalar stabilizer.

Let $W_B$ be the coordinate projection of $W$. The identities
$\langle w,1_Bw'\rangle=0$ and $\langle w,1_Bw'w''\rangle=0$ follow because $1_B$ preserves $V$. Hence $W_B$ satisfies the same quadratic and cubic identities. If it is nonzero, $1_B\notin W_B$, since its squared-coordinate sum is $|B|\ne0$. Mirandola and Z\'emor's theorem applied on $B$ to $W_B$ and $W_B+\C1_B$ gives
\[
 \dim V_B\geq\dim W_B+(\dim W_B+1)-1=2\dim W_B.
\]
On the other hand $V_B\subseteq W_B^\perp$, so $3\dim W_B\leq|B|$. Sum over the blocks, using $\dim W\leq\sum_B\dim W_B$. Blocks with $W_B=0$ require no argument.
\end{proof}

For a level space $L$, let $K$ be its doubly centred kernel. The map taking row and column sums has image in the product of two sum-zero spaces, of dimension $2n-2$. Thus
$\dim L\leq\dim K+2n-2$. The formulas for $S_2,S_3$ and Lemma~\ref{lem:fermat} give $\dim K\leq\lfloor n^2/3\rfloor$. This gives the bound $\dim L\leq\lfloor n^2/3\rfloor+2n-2$; Theorem~\ref{thm:weighted-bound} improves it by keeping the marginal sums within the weighted product argument.
If $L$ is complex, $L\cap\overline L$ is the complexification of its real fixed subspace; that real subspace is level. Therefore $\dim_\C(L\cap\overline L)\leq2n-2$ as well.

The standard hook consists of matrices supported on the first row and first column, with first-row sum zero and the first-column sum below the first row zero. Its dimension is $2n-3$. Its first rook number is zero, its second is the product of the two indicated sums, and all higher ones vanish by support. This proves the lower bound.

\section{The unconditional cubic bound, after Mirandola and Z\'emor}
The coordinatewise-product theorem can be applied before imposing double centring. Keeping the row and column sums as additional weighted coordinates gives a smaller unconditional dimension bound. The weights have both signs, so a stabilizer argument must control their possible zero-sum blocks.

\begin{lemma}[Weighted moments and the product stabilizer]
\label{lem:weighted-moment-stabilizer}
Let $w_1,\ldots,w_M\in\mathbb C^*$ and let $V\subseteq\mathbb C^M$
be a linear space on which
\[
 \sum_{\nu=1}^M w_\nu v_\nu
 =\sum_{\nu=1}^M w_\nu v_\nu^2
 =\sum_{\nu=1}^M w_\nu v_\nu^3=0.
\]
Products of vectors are coordinatewise, and products of subspaces
denote their linear spans. Put
\[
 U=\mathbb C\mathbf1+V,\qquad
 W=U^2,\qquad H=\operatorname{St}(W).
\]
If $H\cap V=0$, then
\[
 3\dim V+\dim H\le M.
\]
In particular $3\dim V\le M-1$.
\end{lemma}
\begin{proof}
Use the nondegenerate bilinear form
$\langle x,y\rangle_w=\sum_\nu w_\nu x_\nu y_\nu$.
Polarization of the three identities gives
$W\subseteq V^{\perp_w}$, hence
\[
 \dim W\le M-\dim V.
\]
The stabilizer $H$ is a unital subalgebra. Since $\mathbf1\in W$,
one has $H\subseteq W$ and $HW=W$.
Set $K=HU$. Then $K$ contains $H+V$, and
\[
 \dim K\ge\dim H+\dim V,\qquad K^2=H^2U^2=W.
\]
The coordinate-product theorem
\citep[Theorem~3.3, p.~7]{MZ} gives
\[
 \dim W=\dim K^2
 \ge2\dim K-\dim\operatorname{St}(K^2)
 =2\dim K-\dim H
 \ge2\dim V+\dim H.
\]
Combining the inequalities proves the assertion. Finally
$\dim H\ge1$ because $\mathbf1\in H$.
\end{proof}

\begin{theorem}[The unconditional cubic bound]
\label{thm:weighted-bound}
For every $n\ge3$, every complex level space of order $n$ satisfies
\[
 3\dim L\le n^2+2n-1.
\]
In particular
$\operatorname{LS}(5)\le11$,
$\operatorname{LS}(6)\le15$, and
$\operatorname{LS}(7)\le20$.
\end{theorem}
\begin{proof}
For $A=(a_{ij})\in L$, put
$r_i=\sum_j a_{ij}$ and $c_j=\sum_i a_{ij}$.
The identity $S_1(A)=0$ gives
$\sum_i r_i=\sum_j c_j=0$.
Define
\[
 b_{ij}=a_{ij}-\frac{r_i+c_j}{2},\qquad
 \alpha=-\frac{n-2}{2n},\qquad
 \lambda=\frac{n^3}{(n-2)^2}.
\]
Map $A$ injectively to
\[
 \Phi(A)=(b_{ij};\,\alpha r_i;\,\alpha c_j)
 \in\mathbb C^M,\qquad M=n^2+2n,
\]
and give the first $n^2$ coordinates weight $1$ and the last $2n$
coordinates weight $-\lambda$.
The map is injective because its last coordinates recover the
marginal sums, after which
$a_{ij}=b_{ij}+(r_i+c_j)/2$.

The quadratic and cubic expansions on $S_1=0$ are
\begin{align*}
 2S_2(A)
 &=\sum_{ij}b_{ij}^2
       -\frac n4\left(\sum_i r_i^2+\sum_j c_j^2\right),\\
 \frac32S_3(A)
 &=\sum_{ij}b_{ij}^3
       +\frac{n-2}{8}\left(\sum_i r_i^3+\sum_j c_j^3\right).
\end{align*}
Here is the cubic calculation explicitly. Put $A_3=\sum_{ij}a_{ij}^3$, $R_3=\sum_i r_i^3+\sum_jc_j^3$, $B=\sum_{ij}a_{ij}r_ic_j$ and $C=\sum_{ij}a_{ij}^2(r_i+c_j)$. Applying the distinct-row and distinct-column partition expansion on three labels gives
\[
 \frac32S_3=A_3-\frac32C+\frac32B+\frac12R_3.
\]
Expanding the cube of $b_{ij}$ gives
\[
 \sum_{ij}b_{ij}^3=A_3-\frac32C+\frac32B+\frac{6-n}{8}R_3.
\]
Indeed $\sum_{ij}a_{ij}(r_i+c_j)^2=R_3+2B$ and $\sum_{ij}(r_i+c_j)^3=nR_3$, since the marginal sums are zero. Their difference proves the displayed cubic identity. Expanding the square gives $\sum b_{ij}^2=\sum a_{ij}^2+(n-4)(\sum r_i^2+\sum c_j^2)/4$, which proves the quadratic identity using \eqref{eq:quad}.

Since $-\lambda\alpha^2=-n/4$ and
$-\lambda\alpha^3=(n-2)/8$, the weighted quadratic and cubic moments
vanish on $V=\Phi(L)$.
Its weighted linear moment vanishes as well: the sum of the
$b$ coordinates and the sum of each marginal coordinate group
are zero. We will also use the unweighted linear constraint
\begin{equation}
 \sum_{\nu=n^2+1}^{M}v_\nu=0\qquad(v\in V).
 \label{eq:weighted-extra-sum}
\end{equation}

Take $U,W,H$ as in Lemma~\ref{lem:weighted-moment-stabilizer}.
It remains to prove $H\cap V=0$.
A unital subalgebra of $\mathbb C^M$ has a basis of disjoint
coordinate projectors, corresponding to a partition into blocks
$B$ \citep[Lemma~2.7, p.~5]{MZ}.
On this basis the restriction of the weighted bilinear form to $H$
is diagonal, with entries
\[
 \omega(B)=\sum_{\nu\in B}w_\nu.
\]
Since $H\subseteq W\subseteq V^{\perp_w}$, every element of
$H\cap V$ belongs to the radical of this restricted form.
It is therefore a linear combination of the indicators of
zero-weight blocks.

If a zero-weight block contains $a$ entry coordinates and $b$
marginal coordinates, then
\[
 a(n-2)^2=bn^3,\qquad a,b>0.
\]
Writing $g=\gcd(n^3,(n-2)^2)$, its counts are positive multiples of
\[
 (a_0,b_0)=
 \left(\frac{n^3}{g},\frac{(n-2)^2}{g}\right).
\]
Here $g=1$ for odd $n$, $g=4$ for $n\equiv0\pmod4$, and $g=8$
for $n\equiv2\pmod4$. For the even cases write $n=2m$ and use $g=4\gcd(2,(m-1)^2)$; for odd $n$, the integers $n$ and $n-2$ are coprime.
The constraint $a_0\le n^2$ rules out such a block except at
$n=4$ or $6$.
At those two orders the minimal counts are respectively
$(16,1)$ and $(27,2)$.
There can therefore be at most one zero-weight block in the
partition.

For every other order the restricted form on $H$ is nondegenerate,
so $H\cap V=0$.
At order four or six its radical is either zero or the line spanned
by the indicator $\mathbf1_B$ of that one block.
But $\mathbf1_B$ has a positive number $b$ of marginal coordinates,
so no nonzero scalar multiple of it satisfies
\eqref{eq:weighted-extra-sum}.
Thus $H\cap V=0$ at these orders too.

Lemma~\ref{lem:weighted-moment-stabilizer} now gives the stronger
intermediate inequality
\[
 3\dim L+\dim H\le n^2+2n,
\]
and $\dim H\ge1$ proves the theorem.
\end{proof}


The conclusion is still quadratic in $n$ and uses no layer beyond $S_3$. At order six it gives the unconditional interval $9\le\LS(6)\le15$. The bound ten in Theorem~\ref{thm:rowcycle} requires a symmetry hypothesis, and no reduction of arbitrary spaces to that class is asserted.


\section{Amplification by direct sums}
Direct sums turn a level-space construction at one order into constructions at all larger orders and determine the limiting dimension ratio.

\begin{proposition}\label{prop:amplify}
If $A\oplus B$ is a block-diagonal matrix, then $R_{A\oplus B}=R_AR_B$. Consequently $\LS(n+m)\geq\LS(n)+\LS(m)$ and the limit in Theorem~\ref{thm:summary} holds.
\end{proposition}
\begin{proof}
A matching splits uniquely between the two diagonal blocks, and its weight is the product of the two matching weights. This proves the generating-function identity. Direct sums of two level spaces are therefore level. In particular, for fixed $r$, $\LS(n)\geq\lfloor n/r\rfloor\LS(r)$, using zero perturbations in the unused rows and columns. Dividing by $n$ and taking the lower limit gives at least $\LS(r)/r$ for every $r$. The upper limit is at most their supremum by definition, proving the equality, including when the supremum is infinite.
\end{proof}

\section{The Taylor-order obstruction and rank-one lines}
This section gives exact tests which prevent limited-degree information from being mistaken for the full level condition.

\begin{proposition}\label{prop:barrier}
Let $q\mid r\leq n$ and let $\omega$ be a primitive $q$th root of unity. Form matrices with $r/q$ independent row blocks
$(a,\omega a,\ldots,\omega^{q-1}a)$ and all remaining rows zero. This is a space of dimension $nr/q$, and $S_k$ vanishes on it unless $q\mid k\leq r$.
Let $C\ge0$. If vanishing in a set of orders $\Omega\subseteq\{1,\ldots,n\}$ forces every subspace to have dimension at most $Cn$, put $a=\lfloor C\rfloor+1$. Then for every $1\leq q\leq\lfloor n/a\rfloor$, $\Omega$ meets $\{q,2q,\ldots,aq\}$. In particular
\[
 |\Omega|\geq\left\lceil\frac{\lfloor n/a\rfloor}{a}\right\rceil,
 \qquad\max\Omega\geq\lfloor n/a\rfloor.
\]
\end{proposition}
\begin{proof}
Cycling the rows within every block multiplies the matrix by $\omega$ or $\omega^{-1}$, depending on the direction. Invariance and homogeneity imply $S_k=\omega^kS_k$, so $q\nmid k$ forces zero. Orders above $r$ vanish because every such submatrix uses a zero row. For the consequence take $r=aq$. The resulting space has dimension $an>Cn$, so at least one tested order must be a positive multiple of $q$ no larger than $aq$. A fixed order can belong to at most $a$ of these sets, since it equals $jq$ for some $1\leq j\leq a$. Counting the sets gives the cardinality bound; the last set gives the maximum-order bound.
\end{proof}

The conclusion is a necessary condition on a proposed method. It neither constructs a superlinear level space nor proves that a suitable growing family of conditions suffices.

\begin{proposition}\label{prop:rankone}
For nonzero $x,y\in\C^n$, the line $\C xy^{\mathsf T}$ is level if and only if $e_k(x)e_k(y)=0$ for every $1\leq k\leq n$, where $e_k$ denotes the elementary symmetric polynomial. Equivalently, the polynomials
\[
 P_x(z)=\prod_i(1+x_i z),\qquad P_y(z)=\prod_j(1+y_jz)
\]
have disjoint nonconstant coefficient supports.
\end{proposition}
\begin{proof}
Each $k\times k$ subpermanent of a rank-one matrix is $k!$ times the product of the selected row factors and selected column factors. Thus $S_k(xy^{\mathsf T})=k!e_k(x)e_k(y)$, and Proposition~\ref{prop:minc} gives the equivalence. Conversely any two nonconstant polynomials with constant term one, degree at most $n$, and disjoint positive coefficient supports factor into such linear factors over $\C$; pad their factor lists with zero entries. This describes every pair, up to coordinate permutation and the rescaling $(x,y)\mapsto(ax,a^{-1}y)$.
\end{proof}
For example, take $x=(1,\omega,\ldots,\omega^{n-1})$ with $\omega$ primitive of order $n$. Its elementary symmetric coefficients vanish except in degrees zero and $n$. Exactly the vectors $y$ with at least one zero coordinate give level lines. Taking $k=n$ in the proposition shows that every rank-one level line has a zero coordinate in at least one factor. Thus two full vectors of roots of unity never give one.

\section{Symmetry and the geometric obstruction}\label{sec:symmetry}
This section identifies the precise group action available at the all-ones matrix.
The group $G=(S_n\times S_n)\rtimes\langle\mathsf T\rangle$ preserves every $S_k$. If an eigenvalue $\omega$ has order $q$, its eigenspace for a coordinate permutation $g$ has one dimension per coordinate orbit of length divisible by $q$. Therefore it has dimension at most $\lfloor n^2/q\rfloor$. If $q>n$, every $S_k$ of positive degree at most $n$ vanishes there by homogeneity; this gives a level space of dimension less than $n$.

This argument includes transposition, but it gives no degeneration of an arbitrary level space. The group $G$ is finite, so every morphism of algebraic groups from a connected torus to $G$ is trivial. A torus in the full orthogonal group of $S_2$ need not preserve $S_3$ or any later condition. Row and column scaling provides the action used for homogeneous permanental Fano schemes by \citet[Section~2.2, Proposition~2.3 and Remark~2.5]{CI}, but it moves the point $J_n$. For a direct example, scaling the first column by two takes the row-supported level line with row $(1,-1,0,\ldots,0)$ to a line with nonzero $S_1$.

\section{The matrix zeta formulation}
This section identifies exactly what the determinant generating function says about the level condition.
For commuting indeterminates $x_1,\ldots,x_n$, put $X=\operatorname{diag}(x_1,\ldots,x_n)$ and
\[
 Z_A(x)=\det(I-XA)^{-1}\in\C[[x_1,\ldots,x_n]].
\]
The inverse exists because the determinant has constant term one.
The squarefree case of MacMahon's master theorem, in the form of
\citet[Corollary~2.2, p.~608]{GJMaster}, states that
\[
 [x_1\cdots x_n]Z_A(x)=\per(A).
\]
Here it has the following precise consequence.

\begin{proposition}\label{prop:zeta}
For a complex linear space $L\subseteq\operatorname{Mat}_n(\C)$, the level condition is equivalent to
\[
 [x_1\cdots x_n]\bigl(Z_{J+t}(x)-Z_J(x)\bigr)=0
 \quad\text{for every }t\in L.
\]
After a linear parametrization of $L$, the homogeneous parts of this equation are exactly
$(n-k)!S_k|_L$, for $1\leq k\leq n$.
\end{proposition}
\begin{proof}
Apply the quoted coefficient identity to $A=J+t$ and $A=J$ and then apply
Proposition~\ref{prop:minc}. A linear parametrization preserves total degree in the parameters, so its degree-$k$ part is precisely the indicated multiple of $S_k|_L$.
\end{proof}

The formulation is therefore an equality of polynomial conditions, including their homogeneous degrees. A bound on the dimension of their common linear subspaces requires a further property of this generating function. The coefficient identity itself gives no additional vanishing orders. For example, at order two its value is
\[
 a+b+c+d+ad+bc
 \quad\text{for }t=\begin{pmatrix}a&b\\c&d\end{pmatrix},
\]
so the required homogeneous equations remain $a+b+c+d=0$ and $ad+bc=0$.

\section{The bound ten under a six-cycle}
This section proves the six-by-six target when the whole level space is invariant under a full cycle of rows or columns. Invariance permits several character components. The proof controls their interactions through the quadratic form and the elementary-symmetric equations.

\begin{theorem}\label{thm:rowcycle}
Let $T$ cyclically permute the six rows of a complex $6\times6$ matrix. If $L$ is a level space and $T(L)=L$, then $\dim L\leq10$. The same statement holds for a full six-cycle of columns.
\end{theorem}

We first prove the small elementary-symmetric assertions on which the theorem rests.
\begin{lemma}\label{lem:singular-e3}
For $m\geq3$, the affine singular locus of $e_3$ in $\C^m$ is the union of the coordinate axes. The cubic hypersurface $e_3=0$ in $\mathbb P^4$ has five ordinary double points and contains no projective plane.
\end{lemma}
\begin{proof}
Put $s=\sum x_i$ and $e=e_2(x)$. The derivative in coordinate $i$ is $e-sx_i+x_i^2$. At a singular point, every coordinate is a root of one quadratic, so there are at most two values. If there is only one value, the derivative is $\binom{m-1}{2}$ times its square, and it is zero only at the origin. If the two values are $a,b$, occurring $r,m-r$ times, their sum and product are $s,e$. Comparing the first and second power sums gives
\[
 (r-1)a+(m-r-1)b=0,
 \qquad (r-1)a^2+(m-r-1)b^2=0.
\]
When both integer coefficients are positive, substitution forces $a=b=0$. If one coefficient is zero, the other value is zero and only one coordinate can be nonzero. Conversely every coordinate axis is singular.

For $m=5$, in the chart at a coordinate point the local equation is $e_2(z_1,\ldots,z_4)+e_3(z_1,\ldots,z_4)$. The quadratic part has matrix $(J_4-I_4)/2$, which is nondegenerate. Thus these five points are ordinary double points.

Suppose a projective plane $\Pi$ were contained in the cubic. In linear coordinates with $\Pi=\{y_1=y_2=0\}$, write the cubic as $y_1Q_1+y_2Q_2$. Its singular points on $\Pi$ are exactly the intersection of the two plane quadrics $q_i=Q_i|_\Pi$. This intersection is finite, since the entire singular locus is finite, so the quadrics form a regular sequence. Their Hilbert series is
\[
 \frac{(1-z^2)^2}{(1-z)^3}=\frac{(1+z)^2}{1-z},
\]
and their zero-dimensional projective intersection has length four. This Hilbert-series identity follows successively from the injective multiplication maps by the two quadrics and needs no genericity assumption.

Each of these intersection points is reduced. At an ordinary double point the four-dimensional tangent space has a nondegenerate quadratic form. The two-dimensional tangent space to $\Pi$ is totally isotropic. Its pairing with a complementary normal space is therefore nonsingular: a vector in the kernel would be orthogonal to both summands and hence zero. This is the Jacobian of $(q_1,q_2)$ along $\Pi$, so each local intersection has length one. It follows that $\Pi$ contains four distinct coordinate points. They span a projective three-space, a contradiction.
\end{proof}

\begin{lemma}\label{lem:evenelementary}
A complex linear subspace of $\C^5$ on which $e_2=e_4=0$ has dimension at most one. A subspace of $\C^6$ with $e_3=e_6=0$ has dimension at most two, and one with $e_2=e_4=e_6=0$ has dimension at most one.
\end{lemma}
\begin{proof}
For the first assertion regard the five coordinates as linear forms on the subspace. If one is identically zero, the equation $e_4=0$ forces a second to be zero, since a product of four nonzero linear forms cannot vanish identically. There remain at most three coordinates; the quadratic form $e_2$ on three coordinates is nondegenerate and has maximal complex isotropic dimension one. The cases of one or two remaining coordinates give the same upper bound directly.

If all five forms are nonzero, divide $e_4=0$ by their product in the fraction field. Grouping proportional forms as $a_iL_C$ gives
\[
 \sum_C\frac{\sum_{i\in C}a_i^{-1}}{L_C}=0.
\]
Multiply by the product of the distinct $L_C$ and reduce modulo each $L_C$. Pairwise relative primeness forces every numerator to be zero. Thus each proportionality class has at least two members, and there are at most two classes. The coordinate forms span the dual of the subspace, so its dimension is at most two. In dimension two the classes have sizes two and three. Write the forms as $aU,-aU,bV,cV,dV$. The reciprocal relation gives $bc+bd+cd=0$. In $e_2$, the cross term vanishes and the $V^2$ coefficient is zero, whereas the $U^2$ coefficient is $-a^2\ne0$. This contradiction proves the first assertion.

On a linear subspace of $\C^6$, the identity $e_6=0$ says that a product of coordinate linear forms is zero in an integral domain. Some coordinate is identically zero. The second assertion is then Lemma~\ref{lem:singular-e3} in the resulting five coordinates, and the third is the first assertion.
\end{proof}

\begin{proof}[Proof of Theorem~\ref{thm:rowcycle}]
Let $\zeta$ be a primitive sixth root of unity and let $v_b=(1,\zeta^b,\ldots,\zeta^{5b})^{\mathsf T}$. The spectral projections of $T$ preserve $L$, so
\[
 L=\bigoplus_{b=0}^5 v_b\otimes U_b
\]
for subspaces $U_b\subseteq\C^6$. Write $d_b=\dim U_b$. Every component is itself level. The rank-one formula of Proposition~\ref{prop:rankone} implies $U_0=0$ and $e_6|_{U_b}=0$ for every $b$. It also gives $e_3=0$ on $U_2,U_4$, and $e_2=e_4=0$ on $U_3$, using
\[
 \prod_{i=0}^5(1+z\zeta^{2i})=(1+z^3)^2,
 \quad\prod_{i=0}^5(1+z(-1)^i)=(1-z^2)^3.
\]
The two lemmas therefore give $d_2,d_4\leq2$ and $d_3\leq1$.

Polarizing $S_2$ between modes one and five gives
\[
 \beta(U_1,U_5)=0,\qquad
 \beta(u,w)=\left(\sum_j u_j\right)\left(\sum_jw_j\right)-\sum_ju_jw_j.
\]
For example the cross coefficient is $-6\beta(u,w)$, obtained by summing $\zeta^i\zeta^{-i'}$ over distinct row indices. The matrix $B=J_6-I_6$ is invertible with $B^{-1}=-I_6+J_6/5$, so $d_1+d_5\leq6$. Equality would imply $U_5=U_1^{\perp_\beta}$. As $e_6|_{U_1}=0$, some coordinate $j$ vanishes throughout $U_1$. Then $B^{-1}e_j\in U_5$, but its entries are $-4/5$ at $j$ and $1/5$ elsewhere; its sixth elementary symmetric function is nonzero. This contradicts $e_6|_{U_5}=0$. Hence $d_1+d_5\leq5$, and
\[
 \dim L\leq5+2+1+2=10.
\]
Transposition proves the column statement.
\end{proof}

\begin{corollary}\label{cor:defect}
For any $6\times6$ level space $L$ and full row six-cycle $T$, put
$r=\dim L-\dim(L\cap TL)$. Then $\dim L\leq10+5r$.
\end{corollary}
\begin{proof}
Intersect $L,TL,\ldots,T^5L$ successively. At each step the dimension drops by at most $r$: the preceding intersection lies in the preceding translate, whose intersection with the next translate has codimension $r$. Their full intersection $K$ has dimension at least $\dim L-5r$, is $T$-invariant, and is level. Apply Theorem~\ref{thm:rowcycle}.
\end{proof}

\begin{proposition}\label{prop:twotransitive}
If a finite two-transitive group on $n\geq2$ row positions preserves a level space $L$, then $\dim L\leq n-1$. The bound is attained by the sum-zero matrices supported in one column.
\end{proposition}
\begin{proof}
The permutation module is $\C\mathbf1\oplus V$, with $V$ the irreducible sum-zero module. Indeed its commuting algebra has dimension two, the number of orbits on ordered pairs, and complete reducibility makes this the sum of the squares of irreducible multiplicities. Therefore any invariant subspace of its tensor product with the column space has the form
$\C\mathbf1\otimes U_0\oplus V\otimes U_1$.
The rank-one conditions force $U_0=0$. For every $2\leq k\leq n$ one can choose a vector of $V$ with nonzero $e_k$: use all $k$th roots of unity in $k$ coordinates and zeros elsewhere. Thus $e_k|_{U_1}=0$ for every $k\geq2$. The polynomial with the coordinates of $u\in U_1$ as roots is $z^n-e_1(u)z^{n-1}$, so $u$ has at most one nonzero coordinate. A linear subspace in a union of coordinate axes lies in one axis, giving $\dim U_1\leq1$ and the result.
\end{proof}

These invariant results do not arise by averaging an arbitrary space. The row-supported level direction $(1,-1,0,\ldots,0)$ has cyclic average with every row $(1,-1,0,\ldots,0)/n$; its second rook number is $-(n-1)/n$, so the average is not level.

\section{Scalar symmetries and maximal hooks}
This section excludes two other ways of producing an eleven-dimensional example.

\begin{theorem}\label{thm:scalar-six}
Let $L$ be a $6\times6$ level space. If some $g\in(S_6\times S_6)\rtimes\langle\mathsf T\rangle$ acts as a scalar $\omega\ne1$ on $L$, then $\dim L\leq10$.
\end{theorem}
\begin{proof}
Let $q$ be the order of $\omega$. For $q\geq4$ the orbit count gives ambient eigenspace dimension at most $\lfloor36/q\rfloor\leq9$. If $q=2$, the $-1$ eigenspace lies in $S_1=0$. Since $g$ preserves its nondegenerate quadratic form $S_2$, that eigenspace is orthogonal to every other eigenspace and its restricted form is nondegenerate. Its dimension is at most eighteen, so its totally isotropic subspace $L$ has dimension at most nine.

Consider $q=3$ without transposition. A row cycle of length $a$ and a column cycle of length $b$ give $\gcd(a,b)$ coordinate cycles of length $\operatorname{lcm}(a,b)$. A primitive cube root therefore has multiplicity
\[
 \sum_{3\mid\operatorname{lcm}(a,b)}\gcd(a,b).
\]
Substitution of the eleven partitions of six shows that multiplicity at least eleven must equal twelve. After exchanging rows and columns, the row permutation has type $3+3$ and the column type is $1^6$, $3+1^3$, or $3+3$. To make this finite check explicit, the maximum multiplicities over all column types for row types
\[
 6,\ 5+1,\ 4+2,\ 4+1+1,\ 3+3,\ 3+2+1,
 \ 3+1+1+1,\ 2+2+2,\ 2+2+1+1,\ 2+1^4,\ 1^6
\]
are respectively $6,4,4,6,12,10,12,6,8,10,12$; the formula also identifies the pairs giving twelve.

In each twelve-dimensional case the restricted cubic $S_3$ has no linear factor. For identity columns, restrict to one row triple and four free columns. The resulting polynomial is $6e_3(a,b,c,d)$. The polynomial $e_3=d(ab+ac+bc)+abc$ is primitive and linear in $d$, with relatively prime coefficients in $\C[a,b,c]$, so it is irreducible. If the column permutation contains a three-cycle, restrict to one row triple and one column triple. The matrix is a circulant $3\times3$ matrix with rows multiplied by $1,\omega,\omega^2$, and its permanent is
\[
 a^3+b^3+c^3+3abc.
\]
This plane cubic is smooth: its gradient equations give $a^2=-bc$, $b^2=-ac$, $c^2=-ab$; multiplying rules out a nonzero solution, and a zero coordinate forces all three to vanish. A reducible plane cubic has intersecting components and is singular, so this polynomial has no linear factor either. A linear factor of the full restricted cubic would restrict to zero, contradicting the nonzero restricted polynomial, or to a linear factor of it. An eleven-dimensional subspace in a twelve-dimensional space is a hyperplane; vanishing there would force just such a factor.

For an element involving transposition, conjugate its coordinate action to $(i,j)\mapsto(j,\pi(i))$. Between distinct cycles of $\pi$ of lengths $a,b$, the two rectangles give $\gcd(a,b)$ cycles of length $2\operatorname{lcm}(a,b)$. Within a cycle square there are $\lfloor a/2\rfloor$ cycles of length $2a$, and one additional cycle of length $a$ if $a$ is odd. Label the cycle by $\Z/a$ to check this: squaring the map translates both coordinates, while their difference is sent to one minus itself. For the eleven cycle types above, the primitive-cube-root multiplicities are
$3,0,0,0,7,4,5,0,0,0,0$. They are at most seven. This completes every case.
\end{proof}

\begin{theorem}\label{thm:hook}
For $p,q\geq2$, let $H_{p,q}$ be the matrices supported in the first row and first column whose first row sums to zero and whose first column below that row sums to zero. If a linear subspace $L$ contains $H_{p,q}$ and $S_1,S_2,S_3$ vanish on $L$ (omitting orders above $\min(p,q)$), then $L=H_{p,q}$. In particular this level space of dimension $p+q-3$ is maximal by inclusion.
\end{theorem}
\begin{proof}
The matching-support argument already proves that $H_{p,q}$ is level. Take $A\in L$ and vary a sum-zero first row $r$. The coefficient of $t$ in $S_2(A+tr)$ is
\[
 \sum_jr_{1j}\sum_{i>1,\ell\ne j}A_{i\ell}.
\]
It follows that $v_j=\sum_{i>1}A_{ij}$ is independent of $j$. Varying a sum-zero column $c$ below the first row likewise makes $R_i=\sum_{j>1}A_{ij}$ independent of $i>1$.
Let $B$ be the block outside the first row and column. When $p,q\geq3$, the coefficient of $tu$ in $S_3(A+tr+uc)$ is
\[
 \sum_{i>1,j>1}c_{i1}r_{1j}T_{ij},
 \quad T_{ij}=\sum_{a>1,\ a\ne i,\ b>1,\ b\ne j}B_{ab}.
\]
The entries $r_{1j}$ for $j>1$ are free, while $c$ ranges through its sum-zero space. Hence $T_{ij}$ is independent of $i$ for each $j$. But
$T_{ij}=\sum B-R_i-v_j+B_{ij}$, so $B_{ij}$ is independent of $i$; equality of its column sums makes every entry one scalar $b$. If $p=2$, this follows from equality of the $v_j$, and if $q=2$ from equality of the $R_i$.

The condition $S_1(A)=0$ now shows that
$A-b(J_{p,q}-pqE_{11})\in H_{p,q}$. Thus $b(J_{p,q}-pqE_{11})\in L$. Counting two-edge matchings gives
\[
 S_2\bigl(b(J_{p,q}-pqE_{11})\bigr)
 =-\tfrac12pq(p-1)(q-1)b^2.
\]
Its coefficient is nonzero, so $b=0$ and $A\in H_{p,q}$.
\end{proof}

Maximality by inclusion does not assert maximum dimension. In particular the theorem rules out extending the standard nine-dimensional hook to dimension eleven, but does not classify all eleven-dimensional subspaces.

\section{Linear symmetries of the level equations}

This section distinguishes preservation of the individual level equations from
preservation of the scheme they define.  The distinction is necessary before a
fixed-point argument can be applied to their Grassmannians.

For a matrix $t=(t_{ij})\in\operatorname{Mat}_n(\mathbb C)$, let
\[
 S_k(t)=\sum_{\substack{R\subseteq\{1,\ldots,n\}\\|R|=k}}
        \ \sum_{\substack{f:R\hookrightarrow\{1,\ldots,n\}}}
             \prod_{i\in R}t_{i,f(i)}.
\]
In particular $S_1$ is the entry sum and $S_n$ is the permanent.

\begin{proposition}\label{prop:conformal-layers}
Let $n\geq3$.  Suppose that $T\in\operatorname{GL}(\operatorname{Mat}_n(\mathbb C))$
and that for each $k$ there is $c_k\in\mathbb C^*$ such that
$S_k(Tt)=c_kS_k(t)$ for all $t$.  Then $T$ is a nonzero scalar times a row
permutation, a column permutation, and possibly transposition.  Conversely every
such transformation has the asserted property.  The scalars satisfy $c_k=c_1^k$.
\end{proposition}

\begin{proof}
At $k=n$ the hypothesis says that $T$ preserves the permanent up to scalar.
Marcus and May's main theorem \cite[Section~2, Theorem, p.~179,
(2.2)--(2.4)]{MarcusMay1962} therefore writes $T(t)$ as $DPtQL$ or $DPt^{\mathsf T}QL$,
where $P,Q$ are permutation matrices and $D,L$ are invertible diagonal matrices.
Undo the permutations and transposition.  The entry $t_{ij}$ is multiplied by
$d_i l_j$.  Comparing its coefficient in $S_1(Tt)=c_1S_1(t)$ gives
$d_i l_j=c_1$ for every pair $i,j$.  Thus all entry multipliers are the same.
The converse and the relation $c_k=c_1^k$ follow from homogeneity.
\end{proof}

The scalar subgroup fixes every vector subspace.  Thus this group supplies no
nontrivial torus action on a Grassmannian.  Preserving the ideal of all equations
is a weaker condition: a generator can change by multiples of lower-degree
generators.  The following finite calculation treats that weaker condition at
order four.

\begin{theorem}\label{thm:scheme-torus-four}
Let $H=\{t\in\operatorname{Mat}_4(\mathbb C):S_1(t)=0\}$ and let
\[
 I=(S_2|_H,S_3|_H,S_4|_H)\subseteq\mathbb C[H].
\]
The Lie algebra of the linear stabilizer of $I$ consists of scalar endomorphisms
of $H$.  Every algebraic torus acting linearly on $H$ and preserving $I$ therefore
acts by scalars, and fixes every point of every Grassmannian of subspaces of $H$.
\end{theorem}

\begin{proof}
Use the fifteen entries other than $t_{44}$ as coordinates
$x_1,\ldots,x_{15}$ on $H$, in row-major order, and put
$t_{44}=-\sum_{i=1}^{15}x_i$.  Write $Q=S_2|_H$ and $C=S_3|_H$.
The homogeneous pieces of $I$ in degrees two and three are
\[
 I_2=\mathbb C Q,\qquad I_3=\mathbb C C+QH^*.
\]
Consequently an infinitesimal endomorphism $A=(a_{ij})$ preserving $I$ must
satisfy
\begin{align}
 \sum_{i,j=1}^{15}a_{ij}x_j\frac{\partial Q}{\partial x_i}
     &=\lambda_2Q,\\
 \sum_{i,j=1}^{15}a_{ij}x_j\frac{\partial C}{\partial x_i}
     &=\lambda_3C+\left(\sum_{j=1}^{15}b_jx_j\right)Q
                                                        \label{eq:lie-cubic}.
\end{align}
These equations are linear in the $242$ unknowns
\[
 (a_{11},a_{12},\ldots,a_{15,15},\lambda_2,\lambda_3,b_1,\ldots,b_{15}).
\]
Here and below the order on the $a_{ij}$ is row-major.

We give the exact matrix recipe for the calculation.  Expand $Q$ and $C$ by the
matching formula above after the indicated substitution for $t_{44}$.  They have
$72$ and $222$ nonzero monomials, respectively.  In the column for $a_{ij}$ put
the coefficients of
\[
 x_j\partial_iQ+x_j\partial_iC.
\]
In the next two columns put the coefficients of $-Q$ and $-C$, and in the last
fifteen columns the coefficients of $-x_jQ$.  Rows are indexed by all monomial
exponent vectors that occur, in lexicographic order; their different total
degrees keep the quadratic and cubic equations separate.  The resulting integer
matrix $M$ has size $780\times242$.

The program \emph{Level-scheme stabilizer}, available through the workbench,
constructs precisely this matrix and records original row and column indices of
a $241\times241$ minor.  Its determinant is $681$ modulo the prime $1009$.
Thus $\operatorname{rank}_{\mathbb C}M\geq241$.  This is an exact integer-matrix
certificate: a nonzero determinant after reduction modulo a prime is nonzero
before reduction.  On the other hand, homogeneity gives the nonzero solution
\[
 A=\operatorname{id}_H,\qquad\lambda_2=2,\qquad\lambda_3=3,
 \qquad b_1=\cdots=b_{15}=0.
\]
Therefore the kernel has dimension exactly one, and every possible $A$ is a
scalar matrix.  Scalars preserve $I$, so this identifies the Lie algebra rather
than merely bounding it.

Finally let an algebraic torus act linearly and preserve $I$.  Decompose $H$ into
its character spaces.  If two distinct characters occurred, an integral
cocharacter of the torus would distinguish them.  Differentiating that
one-parameter subgroup over $\mathbb C$ would give a nonscalar element of the
Lie algebra just computed.  This is impossible.  All characters are equal, so
the torus acts by scalars.
\end{proof}

The theorem concerns the displayed scheme ideal.  It does not identify the
radical of that ideal, classify nonlinear automorphisms of its Fano schemes, or
show that arbitrary level spaces specialize to eigenspaces.  The diagonal
row-and-column torus used by Chan and Ilten
\cite[Section~2.2, Proposition~2.3 and Remark~2.5]{CI} acts on the
homogeneous permanental Fano scheme.  It moves the all-ones base point and does
not preserve the translated level equations here.

\section{Two restrictions on a possible six-dimensional space at order four}

The following restriction records explicitly where a small-case argument still
needs a classification.  Let $B$ be the polarization of $S_2$ on $H$, so that
$B(t,u)=S_2(t+u)-S_2(t)-S_2(u)$.  It is nondegenerate.

\begin{proposition}\label{prop:normal-rank-four}
Suppose $L\subseteq H$ is a six-dimensional level space of $4\times4$ matrices.
For each row and each column there is a linear functional annihilating its
projection of $L$ and taking value one on the all-ones vector.  Let $M\subseteq H^*$
be the span of any eight such functionals, restricted to $H$.  Identify $H^*$ with
$H$ by $B$.  The Gram form on $M$ has rank at most three.
\end{proposition}

\begin{proof}
A row projection of $L$ cannot contain the all-ones vector: otherwise it contains
its negative, and $J+t$ has a zero row for some $t\in L$, contrary to its permanent
being $24$.  Linear separation from a subspace supplies a functional vanishing
on the projection and nonzero on the all-ones vector; rescale it to value one.
The column argument is identical.

Put $r=\dim M$ and let $\rho$ be the rank of its induced Gram form.  The subspace
$U=M^\perp$ has dimension $15-r$ and radical
$M\cap M^\perp$ of dimension $r-\rho$.  Hence its largest totally isotropic
subspace has dimension
\[
 (r-\rho)+\left\lfloor\frac{15-2r+\rho}{2}\right\rfloor
       =\left\lfloor\frac{15-\rho}{2}\right\rfloor.
\]
The inclusion $L\subseteq U$ and the equality $S_2|_L=0$ give
$6\leq\lfloor(15-\rho)/2\rfloor$, so $\rho\leq3$.
\end{proof}

For explicit equations, put $D=\frac13J_4-I_4$.  On the full matrix space the
inverse of the polarization matrix is $D\otimes D$.  Passing to $H$ subtracts
$1/144$ from the pairing of any two coefficient vectors whose sums are one.
Thus normalized row functionals with coefficient vectors $a_i,a_k$ have pairing
\[
 D_{ik}\left(\frac13-a_i\cdot a_k\right)-\frac1{144}.
\]
The other entries follow from the same tensor formula.  Vanishing of the
$4\times4$ minors of this eight-generator Gram matrix is therefore a concrete
necessary condition on a six-dimensional candidate.  It does not yet classify
such candidates.

There is a separate obstruction to proving the bound only by isolating the
standard five-dimensional hook.  Its Grassmannian tangent space has ambient
dimension $5(15-5)=50$.  Exact differentiation of the layer equations gives ranks
$15$ for $S_2$, $33$ for $S_2,S_3$, and again $33$ after $S_4$.  Thus the tangent
space to the full level-space scheme there has dimension $17$.  The quartic has
zero first derivative at a hook because it would require three hook entries on
distinct rows and columns.  These are exact rational linear-algebra calculations,
with a reproducible coefficient construction in the workbench.  Higher-order
obstructions, rather than first-order isolation, are needed at this point.


\subsection{Why quadratic and projection tests do not settle order four}

The normal-space rank condition does not by itself exclude a six-dimensional
space, even if all forbidden zero rectangles are tested.  Here is an explicit
boundary of that approach.

\begin{proposition}\label{prop:quadratic-projection-boundary}
There is a six-dimensional subspace $L\subseteq\operatorname{Mat}_4(\mathbb C)$
on which $S_1=S_2=0$, such that no matrix $J+t$ with $t\in L$ has an
$r\times s$ zero rectangle with $r+s\geq5$.  The space $L$ is not level.
\end{proposition}

\begin{proof}
Let $\omega$ be a primitive cube root of unity.  Let $H_0$ be the standard
five-dimensional hook: its first row sums to zero, its first column below that
row sums to zero, and all other entries vanish.  Let $B$ have zero first row and
column and lower-right block
\[
 \begin{pmatrix}1\\-1\\0\end{pmatrix}
 \begin{pmatrix}1&\omega&\omega^2\end{pmatrix}.
\]
Put $L=H_0\oplus\mathbb C B$.  The supports make this a direct sum.  The matrix
$B$ is doubly centred and
\[
 \sum_{i,j}B_{ij}^2=(1+1)(1+\omega^2+\omega^4)=0.
\]
Consequently $S_1(B)=S_2(B)=0$.  Its $S_2$ pairing with $H_0$ vanishes because
its row and column sums are zero and the supports are disjoint.  Thus
$S_1=S_2=0$ on $L$.

Write $t=r+c+zB$, with $r$ the first-row part of the hook and $c$ its remaining
first-column part.  In the last three columns, the lower three rows of $J+t$ are
\[
 \begin{pmatrix}
  1+z&1+z\omega&1+z\omega^2\\
  1-z&1-z\omega&1-z\omega^2\\
  1&1&1
 \end{pmatrix}.
\]
The first row sums to four, the three lower entries in the first column sum to
three, and the displayed entries rule out every other zero row or column.
A $2\times3$ rectangle containing the last row has a nonzero entry.  One
containing both middle rows cannot vanish in any nonfirst column.  The remaining
possibility consists of the first row and one middle row; its columns include
two nonfirst columns, whose distinct cube-root coefficients prevent both entries
in that middle row from vanishing.  A $3\times2$ rectangle contains two lower
rows and a nonfirst column, so the same alternatives rule it out.  These are all
minimal rectangles with $r+s=5$, and every larger forbidden rectangle contains
one of them.

Finally take $r=E_{12}-E_{11}$ and $c=E_{21}-E_{41}$.  Then
\[
 t(z)=\begin{pmatrix}
 -1&1&0&0\\
 1&z&z\omega&z\omega^2\\
 0&-z&-z\omega&-z\omega^2\\
 -1&0&0&0
 \end{pmatrix}.
\]
Direct expansion of the matching sums gives
\[
 S_1(t(z))=S_2(t(z))=0,\qquad
 S_3(t(z))=z-2z^2,\qquad S_4(t(z))=2z^2.
\]
For example, the coefficient of $z$ in $S_3$ comes from one entry of each of
$r,c,B$; the coefficient of $z^2$ is the contribution from $r$ and two entries
of $B$.  The only nonzero fourth-order terms use the last row at column one
and the first row at column two, giving $2z^2$.  Since
$\operatorname{per}(J+t)=24+6S_1+2S_2+S_3+S_4$, it follows that
\[
 \operatorname{per}(J+t(z))=24+z.
\]
Thus $L$ is not level.  At $z=-24$ the permanent vanishes despite the absence of
any forbidden zero rectangle.
\end{proof}

The example satisfies the normalized projection-annihilator conditions and the
quadratic Gram-rank condition.  It shows that those conditions detect support
obstructions but do not detect cancellation among permanent terms.  A proof of
the order-four bound must also control the mixed cubic equations.

\subsection{A mixed five-cycle family}

\begin{lemma}
If $n\geq3$ and $U\subseteq\mathbb C^n$ is linear with $e_{n-1}|_U=0$, then $\dim U\leq n-2$.
\end{lemma}
\begin{proof}
Let $\ell_1,\ldots,\ell_n$ be the coordinate linear forms on $U$.
If one is zero, at least two are zero: with exactly one zero, the sole surviving term of $e_{n-1}$ is the nonzero product of the remaining forms. Thus $\dim U\leq n-2$ in this case.
Otherwise partition the forms into proportionality classes, writing $\ell_i=a_iL_C$ for $i\in C$, where $a_i\ne0$. In the fraction field of $\mathbb C[U]$,
\[
 0=\frac{e_{n-1}(\ell)}{\ell_1\cdots\ell_n}
   =\sum_C\frac{\sum_{i\in C}a_i^{-1}}{L_C}.
\]
Multiply by $\prod_C L_C$ and reduce modulo each $L_C$.
Distinct representatives are relatively prime, so $\sum_{i\in C}a_i^{-1}=0$ for every class.
Each class has at least two members. Since the coordinate forms span $U^*$, we obtain $\dim U\leq\lfloor n/2\rfloor\leq n-2$.
\end{proof}

\begin{proposition}
Let $\omega$ be a primitive fifth root of unity. Every level space of six-by-six matrices whose rows have the form
\[
 c,\omega c,\omega^2c,\omega^3c,\omega^4c,d,
 \qquad c,d\in\mathbb C^6,
\]
has dimension at most nine.
\end{proposition}
\begin{proof}
The rank-one formula and $\prod_{a=0}^4(1+z\omega^a)=1+z^5$ give
\[
 S_1=\sum_jd_j,\qquad S_5=5!e_5(c).
\]
For the second identity, terms involving the last row and four of the first five rows have coefficient $e_4(1,\omega,\ldots,\omega^4)=0$.
The projection of the level space to $c$ therefore has dimension at most four by the lemma.
Its kernel consists of last-row vectors of sum zero and has dimension at most five. Rank-nullity proves the bound.
\end{proof}

\subsection{A small-characteristic obstruction to specialization certificates}

Consider the closed subscheme $Z_d\subseteq\operatorname{Gr}(d,36)_{\mathbb Z}$ defined by vanishing of the restrictions of the integer polynomials $S_1,\ldots,S_6$ to the universal subspace.
The morphism $Z_d\to\operatorname{Spec}\mathbb Z$ is projective, hence proper.
If $Z_d$ has a complex point, its image contains the generic point of $\operatorname{Spec}\mathbb Z$; its closed image is then all of $\operatorname{Spec}\mathbb Z$.
Thus emptiness over $\overline{\mathbb F}_p$ at any one prime would prove emptiness over $\mathbb C$.
For $d=11$, the primes two, three and five cannot be used in this way.

\begin{proposition}
Over each of $\mathbb F_2,\mathbb F_3,\mathbb F_5$, the equations $S_1=\cdots=S_6=0$ contain an eleven-dimensional linear subspace of six-by-six matrices.
\end{proposition}
\begin{proof}
In characteristic $p$, a block of $p$ identical rows contributes zero to every rook sum involving a nonempty selection from that block.
Indeed cyclic rotation of its row positions acts freely on the corresponding matchings and leaves each monomial value unchanged.
A nonempty matching cannot be fixed by a nonidentity rotation: invariance would assign the same column to all rows in a nontrivial row orbit, contradicting distinctness of columns.
Each orbit therefore contributes $p$ identical terms, hence zero.

In characteristic two, three independently chosen identical-row pairs give a space of dimension eighteen on which every positive rook sum vanishes.
In characteristic three, two independently chosen identical-row triples give dimension twelve.
In characteristic five, choose the first five rows equal and arbitrary, and choose the last row arbitrarily subject to its coordinate sum being zero. This gives dimension $6+5=11$.
The terms using the five-row block vanish, while the only possible nonempty matching confined to the last row contributes its zero first rook sum.
Taking an eleven-dimensional subspace in the first two constructions completes the proof.
\end{proof}

These are vanishing statements for the individual integer-defined layers. They do not use the equality $\operatorname{perm}(J+t)=6!$, which degenerates in these characteristics. An emptiness certificate by this specialization method must use a prime at least seven; no such certificate is proved here.

\subsection{Exploratory unrestricted numerical search}

An original numerical search was carried out on complex isotropic Grassmannian charts that impose $S_1=S_2=0$ up to floating-point roundoff. On the $35$-dimensional hyperplane $S_1=0$, the quadratic matrix is $Q=I-R^{\mathsf T}R-C^{\mathsf T}C$, where $R,C$ take row and column sums. After reducing this nondegenerate form to the identity, a Witt chart has basis
\[
 T=E+F\left(K-\tfrac12B^{\mathsf T}B\right)+HB,
 \qquad K^{\mathsf T}=-K.
\]
Here the columns of $E$ and $F$ are dual eleven-dimensional isotropic frames, and the columns of $H$ form an orthonormal frame for their thirteen-dimensional orthogonal complement. The matrix $K$ is eleven-by-eleven and skew-symmetric, and $B$ is thirteen-by-eleven. For dimension eleven, this gives $55+143=198$ complex parameters.
The implementation evaluates the permanent at seven roots of unity by Ryser's subset formula and extracts $S_3,\ldots,S_6$ by Fourier inversion, with analytic derivatives.
Each of sixteen starts used $1024$ training parameter vectors and an independent holdout of $512$ vectors after Hermitian orthonormalization of the final basis. Eight initial charts extended the known nine-dimensional hook by two isotropic directions, and eight used random charts; all chart parameters were allowed to vary.

No candidate was found. The smallest Euclidean norm of the four holdout root-mean-square layer residuals was $1.869$ for the hook-based starts and $4.206$ for the random starts. By comparison, a dimension-nine calibration with $512$ training points attained $1.65\times10^{-11}$. An earlier calibration with only $100$ training points had a tiny training residual but a holdout residual $4.14\times10^{-4}$; this detected overfitting is documented with the calculation.
The numerical kernel was independently compared with direct enumeration of all square subpermanents, with maximum absolute error $6.35\times10^{-12}$ on the check instances, and its analytic derivative was checked by centered finite differences.
The searches are discovery calculations, not infeasibility certificates: they neither cover the Grassmannian nor rule out other critical points or exact level spaces. Reproducible code, seeds, all start records and open CSV basis data accompany the paper.

\section{Hypermatching subdivision and abelian row symmetries}
This section expresses phase-block substitution as a hypermatching identity and obtains dimension bounds from abelian row symmetries.
For an $m\times n$ array $c$ and an integer $q\geq1$, let
$h_k^{(q)}(c)$ sum $\prod_I\prod_{j\in C_I}c_{Ij}$ over $k$
distinct rows $I$ and pairwise disjoint $q$-subsets $C_I$ of the
columns. Put $h_0^{(q)}=1$.
\begin{proposition}
Let $\omega$ be a primitive $q$th root of unity and set
$t_{(I,a),j}=\omega^ac_{Ij}$ for $0\leq a<q$.
Then
\[
 S_\ell(t)=0\quad(q\nmid\ell),
\]\[
 S_{qk}(t)=(q!)^k(-1)^{k(q-1)}h_k^{(q)}(c).
\]
\end{proposition}
\begin{proof}
For any selected $r$ rows in one phase block, summing their phase
products gives
\[e_r(1,\omega,\ldots,\omega^{q-1}).\]
This vanishes
for $0<r<q$, and equals $(-1)^{q-1}$ at $r=q$. Each full block
has $q!$ bijections to its assigned $q$ columns.
\end{proof}
If $q=ab$, define $(\iota c)_{(I,r),j}=\omega^rc_{Ij}$ for
$0\leq r<b$. The same physical matrix, after a row permutation,
is an $a$-phase matrix with root $\omega^b$. Hence
\[
 h_k^{(a)}(\iota c)=0\quad(b\nmid k),
\]\[
 h_{bj}^{(a)}(\iota c)
 =(-1)^{j(b-1)}
 \left(\frac{(ab)!}{(a!)^b}\right)^j h_j^{(ab)}(c).
\]
This identity increases the number of array rows by $b$ and forces
correlations between them. It supplies no dimension induction
without a bound retaining these correlations.

\begin{theorem}\label{thm:abelian-rows}
Let $n\geq3$ and let $K\subseteq\operatorname{Mat}_n(\mathbb C)$
be a linear level space invariant under the regular row action of a
finite abelian group $G$ of order $n$. If $t$ is the number of
nontrivial characters of $G$ of order two, then
\[
 \dim K\leq\frac{n-1-t}{2}(n-1)
       +t\left\lfloor\frac{n-1}{2}\right\rfloor.
\]
In particular, invariance under a full row cycle gives
\[
 \dim K\leq\left\lfloor\frac{(n-1)^2}{2}\right\rfloor.
\]
For any level space $L$, if $T$ is a full row cycle and
$r=\dim L-\dim(L\cap TL)$, then
\[
 \dim L\leq\left\lfloor\frac{(n-1)^2}{2}\right\rfloor+(n-1)r.
\]
\end{theorem}
\begin{proof}
Index the rows by $G$. A character is a homomorphism $\chi:G\to\mathbb C^*$, and $\chi\otimes u$ denotes the matrix with entry $\chi(g)u_j$ in row $g$ and column $j$. Character projections preserve $K$, so
\[
 K=\bigoplus_{\chi\in\widehat G}\chi\otimes U_\chi,
 \qquad U_\chi=\{u\in\mathbb C^n:\chi\otimes u\in K\}.
\]
The trivial mode is zero: a constant-row matrix has
$S_k=\binom nk k!e_k(u)$, and all elementary symmetric functions
vanishing forces $u=0$. For each nontrivial mode,
$S_n(\chi\otimes u)$ is a nonzero scalar times $\prod_j u_j$.
Since the coordinate ring of $U_\chi$ is an integral domain,
$U_\chi$ lies in some coordinate hyperplane $H_i=\{u:u_i=0\}$.

The polarization of $S_2$ pairs only opposite nontrivial characters,
and its pairing there is $nQ(u,v)$, where
$Q(u,v)=u\cdot v-(\sum_j u_j)(\sum_j v_j)$ has matrix $Q=I-J_n$.
We have $Q^{-1}=I-J/(n-1)$, so every entry of $Q^{-1}e_j$ is
nonzero for $n\geq3$. If $U_\chi\subseteq H_i$ and
$U_{\chi^{-1}}\subseteq H_j$, then
$U_{\chi^{-1}}\subseteq U_\chi^{\perp_Q}\cap H_j$.
The hyperplane $H_j$ cannot contain $U_\chi^{\perp_Q}$, since that
would put $Q^{-1}e_j$ in $U_\chi\subseteq H_i$.
Thus $\dim U_\chi+\dim U_{\chi^{-1}}\leq n-1$.
The conclusion also holds when either space is zero.

For each nontrivial self-inverse character, $U_\chi$ is
totally isotropic in $H_i$. The restriction of $Q$ to $H_i$ is
$I_{n-1}-J_{n-1}$, which is nondegenerate for $n\geq3$, giving
$\dim U_\chi\leq\lfloor(n-1)/2\rfloor$.
There are $(n-1-t)/2$ opposite pairs and $t$ self-inverse
nontrivial characters. Summation gives the general bound. A cyclic
group has $t=0$ for odd $n$ and $t=1$ for even $n$, giving the stated
cycle bound.

For general $L$, put $K=\bigcap_{i=0}^{n-1}T^iL$.
Each successive intersection loses at most $r$ dimensions:
the current intersection lies in $T^jL$, and the next cut is by
$T^{j+1}L$. Therefore $\dim K\geq\dim L-(n-1)r$.
The space $K$ is an invariant level space, so the first bound applies.
\end{proof}
For $n=6$ this gives $12+5r$, weaker than the separately established
$10+5r$ estimate. No improvement to the unrestricted dimension bound is
claimed here.

\section{The Klein quartic in an eight-by-eight slice}
\label{sec:klein}
This section identifies the permanent level conditions on a conjugate pair of three-dimensional representations of $\operatorname{PSL}_2(\mathbb F_7)$. Klein's invariant ring supplies the geometry; the permanent calculations determine which of its invariants occur.

Put $s=\sqrt{-7}$ and $\alpha=(-1+s)/2$. Index rows and columns by
$P=\mathbb F_7\cup\{\infty\}$, in that order, and let
$\rho(g)T=P_gTP_g^{\mathsf T}$. Let $u:z\mapsto z+1$ and
$v:z\mapsto-1/z$, and let $C$ be the conjugacy class of $u$. Define
\[
 D=\sum_{g\in C}\rho(g)-\sum_{g\in C}\rho(g^{-1}),\qquad
 W=\ker(D-8sI).
\]
The character table in Elkies~\cite[\S1.1, p.~53, (1.1)]{Elkies1999}
shows that $W$ is one of the two non-real irreducible summands of dimension three in $\operatorname{Mat}_8(\mathbb C)$.

Here are coordinates over $\mathbb Q(s)$, so the normalization of every coefficient below is specified. In permutation-tuple notation put
$a=(1,0,4,7,2,6,5,3)$, let $h:z\mapsto4z$, and set
\[
 B_0=\frac18(I+\rho(a))(D^2+8sD)E_{02},\quad
 B_1=\rho(h)B_0,\quad B_2=\rho(h^2)B_0.
\]
Write $B(x)=x_0B_0+x_1B_1+x_2B_2$ and
\begin{align*}
 F(x)&=x_0^4+x_1^4+x_2^4
       +3\alpha(x_0^2x_1^2+x_0^2x_2^2+x_1^2x_2^2),\\
 H(x)&=\det\left(\frac{\partial^2F}{\partial x_i\partial x_j}\right)_{0\leq i,j\leq2},\\
 c&=2688(31+3s),\qquad d=-\frac{832}{27}(87+91s).
\end{align*}
A bar on a polynomial or matrix conjugates its coefficients, leaving its variables fixed.

\begin{theorem}\label{thm:klein-pure}
The matrices $B_0,B_1,B_2$ are a basis of $W$, and
\[
 S_1|_W=S_2|_W=S_3|_W=S_5|_W=S_7|_W=0,
 \qquad S_4(B(x))=cF(x),\qquad S_6(B(x))=dH(x),
\]
\[
 S_8|_W=\frac5{1764}(S_4|_W)^2.
\]
The projective level locus in $\mathbb P(W)$ consists of exactly the twenty-four ordinary flexes of Klein's quartic $F=0$. In particular no two distinct level lines in $W$ span a level plane. The same assertions hold on $\overline W$.
\end{theorem}

\begin{proof}
The conjugacy classes are $1A$, $2A$, $3A$, $4A$, $7A$ and $7B$.
The permutation character on $P$ has values $(8,0,2,0,1,1)$ in this order. Its square has inner product
$(192-24)/168=1$ with each three-dimensional character. The central operator
$D$ acts on these summands by $8s$ and $-8s$, respectively, and by zero on all other irreducibles. Thus $W$ has dimension three. The symmetric square of the permutation representation has character $(36,4,3,0,1,1)$; its inner product with either three-dimensional character is
$(108-84-24)/168=0$. Hence the matrices in $W$ are skew-symmetric. Their row and column sums vanish because those sums define equivariant maps to the permutation representation, which has no three-dimensional constituent.

Direct evaluation of the displayed class-sum expression gives
\[
 \sum_{r,j}(B_i)_{rj}(B_k)_{rj}=0,
 \qquad
 \sum_{r,j}(B_i)_{rj}(\overline B_k)_{rj}=1792\delta_{ik}.
\]
In particular the three matrices are independent. Their generator matrices are
\[
 \rho(u)|_W=
 \begin{pmatrix}
 \alpha/2&\alpha/2&0\\
 1/2&-1/2&\overline\alpha/2\\
 1/2&-1/2&-\overline\alpha/2
 \end{pmatrix},\qquad
 \rho(v)|_W=
 \begin{pmatrix}0&0&-1\\0&-1&0\\-1&0&0\end{pmatrix}.
\]
The quartic $F$ is invariant under these matrices. Klein's invariant ring,
as described by Elkies~\cite[\S1.2, pp.~55--58, (1.9)--(1.18)]{Elkies1999},
has generators of degrees $4,6,14,21$. In degrees $1$ through $8$ its only nonzero homogeneous pieces have degrees $4,6,8$; they are one-dimensional, generated by $F,H,F^2$, respectively. Each $S_k|_W$ is invariant. It therefore remains only to determine three constants.

For clarity, the calculation can be performed by the following exact recurrence, with coefficients in $\mathbb Z[s]$ and $s^2=-7$. Put
$P_{0,\varnothing}(t)=1$, and, for $|A|=r$, put
\[
 P_{r,A}(t)=\sum_{j\in A}
       \bigl(1+t(B_0)_{r-1,j}\bigr)P_{r-1,A\setminus\{j\}}(t).
\]
Expanding along the last row proves that
$P_{8,P}(t)=\operatorname{perm}(J+tB_0)$. Dividing its coefficient of $t^k$ by $(8-k)!$ gives
\[
\begin{array}{c|l}
 k&S_k(B_0)\\\hline
 4&83328+8064s\\
 6&-2502656+2396160s\\
 8&18391040+3809280s.
\end{array}
\]
Since $F(1,0,0)=1$ and $H(1,0,0)=-216(3+s)$, these values give the stated $c,d$ and $5/1764$. The same values follow from the inclusion-exclusion formula for the permanent.

Consequently a nonzero $B(x)$ is a level direction precisely when
$F(x)=H(x)=0$. The quartic is smooth, and its intersection with its Hessian consists of its twenty-four ordinary flexes
\cite[\S2.1, Proposition (i), pp.~64--65]{Elkies1999}; see also
\cite[Lemma 1.4, p.~32]{Kaneta1998}.
A level plane would give a projective line contained in this smooth irreducible quartic, which is impossible. Conjugation gives the other summand.
\end{proof}

The invariant ideal of these level conditions has a particularly short description. With Klein's generators $\Phi_4,\Phi_6,\Phi_{14},\Phi_{21}$,
\[
 I=(\Phi_4,\Phi_6)\subset\mathbb C[W]^G,
 \qquad
 \mathbb C[W]^G/I\simeq
 \mathbb C[\Phi_{14},\Phi_{21}]/(\Phi_{21}^2-\Phi_{14}^3).
\]
The last relation is Elkies's equation (1.18) modulo $(\Phi_4,\Phi_6)$.
Thus the invariant ideal is prime, although its extension to the polynomial ring cuts out the cone over twenty-four distinct points.

\begin{proposition}\label{prop:klein-mixed}
For $T=B(x)+\overline B(y)$, where $x,y\in\mathbb C^3$ are independent coordinates, put $q(x,y)=\sum_i x_i y_i$. Then
\[
 S_2(T)=1792q(x,y),
\]
\[
 S_4(T)=cF(x)+\overline c\,\overline F(y)+774144q(x,y)^2.
\]
Every level subspace of $W\oplus\overline W$ has dimension at most two.
\end{proposition}

\begin{proof}
The quadratic formula follows by polarizing
$S_2(T)=\frac12\sum_{i,j}T_{ij}^2$ on doubly centred matrices and using the basis pairings above. For the quartic, the representation identities
$\operatorname{Sym}^2W\simeq\mathbf6$ and
$\operatorname{Sym}^3W\simeq\mathbf7\oplus W^*$
\cite[p.~55]{Elkies1999}
show that the invariant spaces in bidegrees $(3,1)$ and $(1,3)$ vanish, while the invariant space in bidegree $(2,2)$ is spanned by $q^2$. The coefficient of $q^2$ is determined by
\[
 S_4(B_0+\overline B_0)=940800,
 \qquad c+\overline c=166656;
\]
their difference is $774144$. These evaluations use the preceding recurrence.

The quadratic condition first gives $\dim L\leq3$. Suppose equality holds.
Neither projection of $L$ to a pure summand can have rank at most one, because its kernel would contain a two-dimensional pure level subspace, excluded by Theorem~\ref{thm:klein-pure}.
If the projection to $W$ has rank three, then $L$ is the graph of $y=Ax$.
Isotropy gives $A^{\mathsf T}=-A$, hence $\operatorname{rank}A\leq2$.
The quartic identity becomes $cF(x)+\overline c\,\overline F(Ax)=0$.
It would make $F$ constant along a nonzero direction in $\ker A$.
A homogeneous plane quartic independent of one coordinate is singular at the corresponding point, contrary to the smoothness of $F=0$.
The same argument applies with the two projections reversed.

Both projections therefore have rank two. Their one-dimensional kernels in $L$ are distinct. The equality of the two restricted quartics makes either polynomial constant along both kernel directions. It is consequently a nonzero fourth power of a linear form on $L$. Nonzero follows because the Klein quartic contains no line. Hence the projective line in either projection meets the quartic at a single point with multiplicity four. This would be a hyperflex. Klein's twenty-four flexes are all ordinary, so it has no hyperflex. This contradiction proves the bound.
\end{proof}

\begin{proposition}\label{prop:klein-noextension}
If a level subspace of $W\oplus\overline W$ contains a nonzero pure level direction from either summand, then it has dimension one.
\end{proposition}

\begin{proof}
It suffices to exclude a level plane $L$ containing a line $\ell\subset W$.
The projection to $\overline W$ has rank one. If the other projection had rank two, the quartic identity would express the restriction of $F$ to a projective line as either zero or a fourth power, again giving a line on the quartic or a hyperflex. Both are impossible. Thus
$L=\ell\oplus\ell'$ for a pure level line $\ell'\subset\overline W$, and the two lines must be orthogonal for $q$.

There are only twenty-four possible $\ell'$. The following finite calculation excludes them without approximating their coordinates. Let $\zeta$ be a primitive seventh root of unity, use
$s=\zeta+\zeta^2+\zeta^4-\zeta^3-\zeta^5-\zeta^6$, and reduce at the prime given by $\zeta\mapsto16$ in $\mathbb F_{29}$; then $s\mapsto15$.
By transitivity we may take $\ell$ to be the $\zeta$-eigenline of $\rho(u)$.
Its reduction is $x=(1,16,3)$. Generate the conjugate orbit using the two displayed generator matrices with $s$ replaced by $-s$, starting at the $\zeta^{-1}$-eigenline. There are twenty-four distinct reductions. Precisely two have dot product zero with $x$, and the permanent recurrence gives
\[
\begin{array}{c|c}
 y&S_6\bigl(aB(x)+b\overline B(y)\bigr)\pmod {29}\\\hline
 (1,15,26)&8a^4b^2\\
 (1,19,24)&11a^2b^4.
\end{array}
\]
Every generator matrix is integral at this prime and invertible after reduction. The orbit has twenty-four points in characteristic zero, and its twenty-four distinct reductions therefore represent every point injectively. An orthogonality relation in characteristic zero survives reduction, so the two displayed candidates include every possible partner. Their nonzero sextic reductions rule out a vanishing sextic in characteristic zero. This excludes the plane; conjugation handles the other summand. The exact orbit enumeration and permanent computations are reproduced by the programs in the workbench.
\end{proof}

A mixed level plane transverse to both pure summands is not excluded by these arguments. Its existence is a separate finite-dimensional problem.

\paragraph{Other finite-group slices.}
Let $G\leq S_n\times S_n$ be finite and let $W$ be an isotypic summand for an irreducible representation not isomorphic to its dual. There are no invariant linear or quadratic forms on $W$, hence $S_1|_W=S_2|_W=0$. Put
$R=\mathbb C[W]^G$ and
$I=(S_3|_W,\ldots,S_n|_W)\subset R$.
Then
\[
 I\mathbb C[W]\cap R=I.
\]
Indeed, averaging an expression $f=\sum h_kS_k$ for an invariant $f$ replaces each $h_k$ by its Reynolds average and preserves the equality. Thus the quotient ideal is exactly the contraction of the level ideal. Each degree $k$ with $R_k=0$ contributes no equation; an explicit invariant presentation reduces the remaining equations to finitely many homogeneous components. The Klein calculation above carries this procedure through for one conjugate pair, without asserting a classification for arbitrary finite groups.


\section{What remains open}
The linear bound for $\LS(n)$ and the unrestricted target $\LS(6)\leq10$, the equality $\LS(4)=5$, and the existence of a seven-dimensional level space of $4\times5$ matrices remain unresolved here. The connected formulas provide the exact next equations, but no dimension inequality for their common isotropic subspaces has been proved. A numerical search without a witness cannot replace such an inequality or a certified emptiness calculation.

The Fourier-minor route has a specific restriction. At composite order $n=ab$, with $a,b>1$, the submatrix of the Fourier matrix on rows $\{0,a\}$ and columns $\{0,b\}$ is all ones, and is singular. Thus the all-minors theorem cannot extend to composite order. \citet[Theorem~B and Corollary~2]{Loukaki} and \citet[Theorem~1.2]{CLMP} establish results for principal or specially matched minors under stated arithmetic hypotheses. They do not provide unrestricted nonvanishing for the induction proposed here.

\section*{Method and records}
The mathematics and text were produced by GPT-6 Astra (ChatGPT Pro mode). Exact rational calculations were used to check the connected expansions; they are checks of the formulas, not a classification of all level spaces. Specified statements have separate re-proofs by the same model, with the precise scope recorded at \url{\recordurl}. No human has checked this paper. The program \emph{Rook-cumulant verifier} and its expected output are available through the workbench at \url{\workbenchurl}. This text is dedicated to the public domain under CC0~1.0.

\begin{thebibliography}{99}
\bibitem[Anthropic(2026)]{Anthropic} Anthropic. (2026). \emph{VP $\ne$ VNP over every field}. Preprint, Theorem~1.1. \url{https://hexagonmath.org/2611.00954v1}.
\bibitem[Caragea et~al.(2025)]{CLMP} Caragea, A., Lee, D.~G., Malikiosis, R., \& Pfander, G.~E. (2025). \emph{Chebotarev type results for composite group order. Part~1: Square free natural numbers}. arXiv:2505.24326, Theorem~1.2.
\bibitem[Chan and Ilten(2015)]{CI} Chan, M., \& Ilten, N. (2015). Fano schemes of determinants and permanents. \emph{Algebra \& Number Theory, 9}(3), 629--679. arXiv:1312.2577v2, Section~2.2, pp.~7--8.
\bibitem[Elkies(1999)]{Elkies1999} Elkies, N.~D. (1999). The Klein quartic in number theory. In S. Levy (Ed.), \emph{The eightfold way: The beauty of Klein's quartic curve} (pp.~51--102). Cambridge University Press. Section~1.2, equations (1.9)--(1.18); Section~2.1, Proposition~(i), pp.~64--65. \url{https://library.slmath.org/books/Book35/files/elkies.pdf}.
\bibitem[Goldman et~al.(1976)]{GJW} Goldman, J.~R., Joichi, J.~T., \& White, D.~E. (1976). Rook theory.~V. Rook polynomials, M\"obius inversion and the umbral calculus. \emph{Journal of Combinatorial Theory, Series A, 21}, 230--239. \url{https://doi.org/10.1016/0097-3165(76)90067-4}.
\bibitem[Goulden and Jackson(1992)]{GJMaster} Goulden, I. P., \& Jackson, D. M. (1992). Immanants, Schur functions, and the MacMahon master theorem. \emph{Proceedings of the American Mathematical Society, 115}(3), 605--612. Corollary~2.2, p.~608.
\bibitem[Kaneta(1998)]{Kaneta1998} Kaneta, H. (1998). A 72-arc associated with the $A_6$-invariant sextic. \emph{RIMS K\^oky\^uroku, 1063}, 30--40. Lemma~1.4, p.~32.
\bibitem[Landsberg et~al.(2013)]{LMR} Landsberg, J.~M., Manivel, L., \& Ressayre, N. (2013). Hypersurfaces with degenerate duals and the geometric complexity theory program. \emph{Commentarii Mathematici Helvetici, 88}(2), 469--484. \url{https://doi.org/10.4171/CMH/292}. Locator: April~2010 author manuscript, Theorem~1.0.1, p.~1. \url{https://math.univ-lyon1.fr/~ressayre/PDFs/lmr.pdf}.
\bibitem[Loukaki(2025)]{Loukaki} Loukaki, M. (2025). Chebotarev's theorem for cyclic groups of order $pq$ and an uncertainty principle. \emph{Bulletin of the London Mathematical Society, 57}(12), 3841--3856. \url{https://doi.org/10.1112/blms.70192}.
\bibitem[Marcus and May(1962)]{MarcusMay1962} Marcus, M., \& May, F. C. (1962). The permanent function. \emph{Canadian Journal of Mathematics, 14}, 177--189. \url{https://doi.org/10.4153/CJM-1962-013-4}.
\bibitem[Mignon and Ressayre(2004)]{MR} Mignon, T., \& Ressayre, N. (2004). A quadratic bound for the determinant and permanent problem. \emph{International Mathematics Research Notices, 2004}(79), 4241--4253. Theorem~1.
\bibitem[Minc(1978)]{Minc} Minc, H. (1978). \emph{Permanents}. Addison--Wesley. The expansion used here is also stated and proved in the public module cited below.
\bibitem[Mirandola and Z\'emor(2015)]{MZ} Mirandola, D., \& Z\'emor, G. (2015). Critical pairs for the Product Singleton Bound. \emph{IEEE Transactions on Information Theory, 61}(9), 4928--4937. Theorem~3.3; arXiv:1501.06419v2, p.~7.
\bibitem[OpenAI(2026a)]{OpenAICubic} OpenAI. (2026a, September 24). \emph{A cubic lower bound for border determinantal complexity of the permanent}. Preprint. Theorem~1.1, p.~2. \url{https://github.com/openai/math/tree/main/preprints/A-cubic-lower-bound-for-border-determinantal-complexity-of-the-permanent-September-24-2026}.
\bibitem[OpenAI(2026b)]{OpenAI} OpenAI. (2026b). \emph{VP and VNP are different}. November~23 edition, Theorem~1.1 and Corollary~1.2. \url{https://github.com/openai/math/tree/main/preprints/VP-and-VNP-are-different-November-23-2026}.
\bibitem[P versus NP programme(2026a)]{Fleet} P versus NP programme. (2026a). \emph{Public proof graph and Lean repository}. \url{https://pequalnp.com}; \url{https://github.com/brilliant-hq/pequalnp-proofs}.
\bibitem[P versus NP programme(2026b)]{FleetTransfer} P versus NP programme. (2026b). Level-space transfer theorem. \emph{Lean repository}, module \texttt{PEqualNP.Gct.L077}, theorem \texttt{PEqualNP.Gct.L077}.
\bibitem[P versus NP programme(2026c)]{FleetExpansion} P versus NP programme. (2026c). Minc expansion. \emph{Lean repository}, module \texttt{PEqualNP.Gct.L084}, theorem \texttt{PEqualNP.Gct.L084}.
\end{thebibliography}
\end{document}
